\section{Major Contributions}\label{sec:contrib}

This project contributes four things, described in full in \cref{sec:design}:
\begin{enumerate}
\item \textbf{A lightweight runtime SpMV selector.} One GPU kernel extracts structural features; a shallow decision
  tree, trained offline and compiled into the library as plain C++, chooses among \stat{p2-configs} configurations of
  CSR, ELL, SELL, ArgCSR, cuSPARSE and two new kernels (ArgCSR with a chunk size per row region, and CSR with lanes
  per row); a break-even test converts formats only when the expected number of calls repays the conversion.
\item \textbf{An overlapped, self-calibrating out-of-core SpGEMM.} Paper~1's chunking, extended with double and
  triple buffering, a cost model that simulates the copy/compute pipeline from bandwidths and kernel rates measured at
  run time, and re-partitioning while it runs, all within a hard device-memory budget.
\item \textbf{A re-evaluation of both papers} on a PCIe-attached consumer GPU against the 2026 cuSPARSE.
\item \textbf{Sequential and OpenMP CPU versions} of the same kernels, for same-machine comparisons.
\end{enumerate}

\subsection{Contributions of the papers we studied}\label{sec:papers}
\paragraph{Paper 1: Deveci et al., multilevel-memory SpGEMM~\cite{deveci2018multilevel}.} The paper runs the
KokkosKernels \texttt{kkmem} SpGEMM~\cite{deveci2018kkmem} on a KNL and on a P100 connected to its host over NVLink.
It contributes (i) two chunked out-of-core algorithms: \emph{chunk1} keeps a row block of $A$ and $C$ resident and
streams the row blocks of $B$, \emph{chunk2} keeps a row block of $B$ resident and streams $A$ and $C$; (ii) a
partitioning heuristic, Algorithm~4, that splits fast memory 75\%/25\% between the resident and the streamed operand
and picks the variant with the smaller copy volume; and (iii) a placement study (its Table~3) showing that $B$ in host
memory costs 7--29\x, while $A$ and $C$ hurt only when large, and that unified memory reaches at most 30\% of the
device's speed.

\paragraph{Paper 2: Heller and Oberhuber, adaptive row-grouped CSR~\cite{heller2012argcsr}.} The paper contributes
the ArgCSR format: rows are packed into groups handled by one CUDA block, spare threads go to the rows with the
longest chunks, and chunks are padded and interleaved so that reads coalesce. Measured on 1,600 matrices on a Fermi
C2070, ArgCSR beat CPU CSR on 1,168 of them and was up to ten times faster than the cuSPARSE of 2012 on circuit
matrices. The paper also documents how strongly its chunk-size parameter matters, and recommends testing several
formats per matrix without automating that choice.

\paragraph{Other work we drew on.} CSR-vector and ELL kernels~\cite{bell2009spmv}, SELL-$C$-$\sigma$~\cite{kreutzer2014sell},
CSR5~\cite{liu2015csr5} and merge-based SpMV~\cite{merrill2016merge} define the candidate formats. SMAT~\cite{li2013smat},
decision trees~\cite{sedaghati2015selection}, SVMs~\cite{benatia2016svm} and CNNs~\cite{zhao2018cnn} have all been used to
select formats; unlike them, our selector also predicts the cost of \emph{acting} on its choice. Hash-based SpGEMM
(nsparse~\cite{nagasaka2017nsparse}, spECK~\cite{parger2020speck}) and output-size handling~\cite{liu2014spgemm} inform the
SpGEMM engine.

\subsection{Experiments: test-bed and datasets}\label{sec:testbed}
\paragraph{Hardware and software.} The primary machine is a laptop with an RTX~4050 Laptop GPU (Ada, sm\_89, 6\,GB,
20~SMs, 24\,MB L2) on PCIe~4.0~x8 and an 8-core, 16-thread Ryzen~7 7840HS, running Ubuntu~22.04 under WSL2 with CUDA~13.3,
cuSPARSE~12.8 and g++~11.4. \Cref{tab:platform} gives its measured characteristics. A Kaggle Tesla~T4 (CUDA~12.8, 4-thread
Xeon) repeats the same-machine comparison, and the OpenMP scaling study also ran on a Core i5-1235U. No simulator is
used: everything is measured on real hardware.

\paragraph{Datasets.} SpMV: \stat{p2-matrices} SuiteSparse matrices~\cite{davis2011suitesparse} for training (50 drawn at
random, real and square with $10^4$--$5\cdot10^6$ nonzeros, plus 10 named in paper~2) and \stat{p3-heldout-matrices}
different matrices drawn the same way for evaluation, never seen in training. CG: the SPD matrices of these sets, 20 more
SPD matrices, and 3-D Laplacians. SpGEMM: paper~1's multigrid products on the 7-point Laplace3D with $2\times2\times2$
aggregation ($P$, $R=P^T$), generated in memory, which allows multi-GB problems.

\paragraph{Procedure.} Kernels are timed with CUDA events: 5 warm-up runs, then the median of 20 with the
interquartile range; the L2 cache is flushed before every timed SpMV. Out-of-core SpGEMM is timed from pinned inputs to
pinned output, including every transfer and the planning time. Methods are run in a freshly shuffled order in every
repetition, because a fixed order produced fake differences of $\pm9\%$ under WSL2. Every result is checked against a
CPU reference, and every CSV records the GPU, driver, library versions and git commit. Runs use AC power and the
``best performance'' mode. \texttt{scripts/run\_all.sh} rebuilds from a clean tree, reruns every test and benchmark,
and regenerates every number, table and figure in this report.

\begin{table}[t]
\centering\small
\caption{Platform measurements of the RTX~4050 laptop (\texttt{results/platform/uvm\_probe.csv}).}\label{tab:platform}
\resizebox{\columnwidth}{!}{%
\begin{tabular}{@{}lr@{}}
\toprule
Device memory, sequential read (4 GiB) & \stat{pf-device-read} GB/s \\
Device memory, copy (read + write) & \stat{pf-device-copy} GB/s \\
Pinned host $\to$ device / device $\to$ host & \stat{pf-h2d} / \stat{pf-d2h} GB/s \\
Both directions at once (one copy engine) & \stat{pf-bidir} GB/s \\
Zero-copy / managed (UVM) read by a kernel & \stat{pf-zc-read} / \stat{pf-uvm-read} GB/s \\
Random 8-byte gathers, device vs zero-copy & \stat{pf-gather-device} vs \stat{pf-gather-zc} M/s \\
\bottomrule
\end{tabular}}
\end{table}
